\section{The nonlinear escaping states in the original two-root quotient}
Retain the original polynomial, all four marked roots and both signs
at each root, as in FS13 \cite{ES488,Tao}. This section evaluates the
actual affine states, rather than a different multiplication operator,
in the original receiver and the complete quotient of OCQ1--11.
The original Gamma matrices and all moments are ECC24--27
\cite{WCF}. Throughout this section the original reference quartet,
its marking $K$, its moments, and its receiving matrix $M_*$ are fixed;
the algebraic ES family ECC44 varies. Its parameter $t$ is complex,
and this family makes no assertion of integral witnesses for arbitrary
$t$. Fix $P\ne0$.

\subsection{Exact finite quotient, complete complement, and both norms}
Put $m_j(r)=v(r)^{\mathsf T}M_*e_j$, $1\le j\le4$, with precisely
$M_*=C_\kappa\Psi_*$ from ECC27. Define
\[
 E_t=\begin{pmatrix}1&P&P^2&P^3\\
 0&1&2P+t&3P^2+3Pt+t^2\end{pmatrix},\quad
 u_j(t)=E_tM_*e_j=
 \begin{pmatrix}m_j(P)\\m'_j(P)+t m''_j(P)/2+t^2m'''_j(P)/6\end{pmatrix}.
 \tag{EQR1}
\]
This is the exact divided difference, including at zero. For either
$b=U,W$ put $\mathcal Q_{b,t}=E_tF_b^{-1}$ on the scalar receiver.
Then $\mathcal Q_{U,t}\Phi_*q=\mathcal Q_{W,t}\Psi_*q=E_tM_*q$.
The equality follows from both receiver identities in ECC27 and
specifies the actual quotient of each original physical vector.

For $r=P$ or $P+t$ and $s^2=A f'(r)$, the image of the complete
original affine state is
\[
 \begin{aligned}
 Y(r,s,t)={}&s^{-1}u_1-r u_2
 +s(-iu_2+2ru_3+7ir^2u_4)\\
 &+s^2\{3iu_3+(B-17r+Ar^2)u_4\}
 +s^3(Au_3-13iu_4)-2As^4u_4.
 \end{aligned}\tag{EQR2}
\]
Every $u_j$ in this formula is evaluated at $t$; $A,B$ are the
literal coefficients on ECC44. Substituting FS13 proves every term.

Here is the full complementary information, including an inverse.
For $f=(f_0,f_1,f_2,f_3)^{\mathsf T}=M_*q$, set
\[
 y=E_tf,\qquad c=f_2+(2P+t)f_3,\quad d=f_3,\qquad
 f(r)=y_0+y_1(r-P)+(r-P)(r-P-t)(c+dr).
 \tag{EQR3}
\]
Coefficient comparison proves both compositions. Its determinant
before $M_*$ is one. Thus no state coordinate has been discarded:
$q=M_*^{-1}[f]$ is the complete inverse. Let $K_t$ have the two
coefficient columns of $(r-P)(r-P-t)$ and $r(r-P)(r-P-t)$,
and let $L_t$ have rows $(0,0,1,2P+t)$ and $(0,0,0,1)$.
Then $E_tK_t=0$, $L_tK_t=I_2$; these are the entire kernel since
$E_t$ has rank two.

Retain $J_{b,t}=E_tH_b^{-1}E_t^*$, $G_{b,t}=J_{b,t}^{-1}$ from
ECC39. Define $R_{b,t}=H_b^{-1}E_t^*G_{b,t}$. Exact multiplication
gives $E_tR_{b,t}=I_2$ and $K_t^*H_bR_{b,t}=0$. For the actual
$f=M_*q$, put $w_b=L_t(f-R_{b,t}y)$. Then
\[
 f=R_{b,t}y+K_tw_b,\quad
 \|F_bf\|_b^2=y^*G_{b,t}y+w_b^*(K_t^*H_bK_t)w_b.
 \tag{EQR4}
\]
The first identity follows because its difference lies in the kernel
and $L_t$ inverts its kernel coordinates; the second is the proved
orthogonality. It retains both complete original norms, every moment,
and every complementary direction. Both displayed Grams are positive
also at zero.

\subsection{Every collision-escape case, with the exact constants}
Write $t=w^2$ on its local double cover and retain both signs of each
square root. For the root $r=P+\delta t$, $\delta\in\{0,1\}$,
the exact ECC44 path has
\[
 s=\lambda w(1+O(w^2)),\qquad
 \lambda^2=(2\delta-1)a,\qquad a=3P/22,
 \quad A_0=-5/(22P),\quad B_0=-3P/2.
 \tag{EQR5}
\]
Indeed the two derivative values are respectively $-tA h(P)$ and
$tA h(P+t)$ and $Ah(P)\to a$. Each nonzero analytic unit has the
two analytic square roots obtained by its convergent power series.
The four choices $(\delta,\lambda)$ are exactly the four signed
collision states. No choice of sign is suppressed.

Let $U_j=u_j(0)=(m_j(P),m'_j(P))^{\mathsf T}$,
$\alpha=m''_1(P)/2$, $\beta=m''_2(P)/2$.
Expanding the exact EQR2 now gives the following exhaustive cases.
If $U_1\ne0$, then
\[
 wY\longrightarrow U_1/\lambda,\qquad
 |t|^{1/2}\|Y\|_{G_{b,t}}
   \longrightarrow \sqrt{U_1^*G_{b,0}U_1}/|\lambda|>0.
 \tag{EQR6}
\]
Only its first term has order $w^{-1}$, which proves the vector
limit and its full norm constant. If $U_1=0$ but $U_2\ne0$, the
first term is $O(w)$ and
\[
 Y\longrightarrow-PU_2,\qquad
 \|Y\|_{G_{b,t}}\longrightarrow |P|\sqrt{U_2^*G_{b,0}U_2}>0.
 \tag{EQR7}
\]
Thus boundedness of this quotient is evaluated rather than assumed.

It remains to calculate the stratum $U_1=U_2=0$. The columns
$M_*e_1,M_*e_2$ are independent and lie in the two-dimensional
kernel of $E_0$, so they form its full basis. Since $M_*$ is
invertible, $U_3,U_4$ are independent. Put
\[
 V_0=2PU_3+7iP^2U_4\ne0,\qquad
 Z_\lambda=\lambda^{-1}\alpha\binom01+\lambda V_0,
 \quad
 W_\lambda=-P\beta\binom01+
 \lambda^2\{3iU_3+(B_0-17P+A_0P^2)U_4\}.
 \tag{EQR8}
\]
Term-by-term expansion of EQR2 proves
\[
 Y=wZ_\lambda+w^2W_\lambda+O(w^3).
 \tag{EQR9}
\]
For detail, $u_1=t\alpha(0,1)^{\mathsf T}+O(t^2)$ gives the
first summand of $Z_\lambda$, $-ru_2$ gives the first summand of
$W_\lambda$, and the terms linear and quadratic in $s$ give their
other summands. Corrections to $r,A,B,s,u_j$ in the linear term
begin at order $w^3$; the cubic and quartic terms do also. These
observations account for every original term through order two.

If $Z_\lambda\ne0$, the exact order is $|t|^{1/2}$, with constant
$\sqrt{Z_\lambda^*G_{b,0}Z_\lambda}$ at receiver $b$.
The equation $Z_\lambda=0$ can hold for at most one root label
$\delta$: multiplication by $\lambda$ gives
$\alpha(0,1)^{\mathsf T}+\lambda^2V_0=0$, and the two possible
values of $\lambda^2$ are opposite and nonzero. If both vanished,
their difference would imply $V_0=0$, a contradiction. Cancellation
at a label occurs for both of its signs, since $Z_{-\lambda}=-Z_\lambda$.

In that cancelling pair, $W_\lambda$ is nonzero. Its first component
can be calculated completely. The first component of $V_0$ is zero,
so $m_3(P)=-7iP m_4(P)/2$. The four values $m_j(P)$ cannot all vanish,
because $v(P)^{\mathsf T}M_*$ is a nonzero row. Thus $m_4(P)\ne0$.
Substitution in EQR8 gives
\[
 (W_\lambda)_1=-\lambda^2\frac{181P}{22}m_4(P)\ne0,
 \quad
 t^{-1}Y\longrightarrow W_\lambda,\quad
 |t|^{-1}\|Y\|_{G_{b,t}}
       \longrightarrow\sqrt{W_\lambda^*G_{b,0}W_\lambda}>0.
 \tag{EQR10}
\]
The scalar follows from $21P/2+B_0-17P+A_0P^2=-181P/22$.
Equations EQR6--10 therefore exhaust the orders
$-1/2,0,1/2,1$ on the specified original ES path and give both
original metric constants. Each alternative is decided by the stated
finite matrices of the fixed original receiver, not by an unspecified
missing assumption.

The location of the bounded stratum is also finite and explicit.
$m_1$ is a nonzero polynomial of degree at most three because
$M_*e_1\ne0$. The equation $U_1=0$ means that $P$ is a repeated
root of that actual polynomial, so there is at most one such $P$.
For completeness, for $m_1(r)=ar^3+br^2+cr+d$ with $a\ne0$,
set $\Delta=b^2-3ac$. When $\Delta\ne0$ the only possible point is
\[
 P=(9ad-bc)/(2\Delta);
 \quad\text{it belongs precisely when }m_1(P)=m'_1(P)=0.
 \tag{EQR11}
\]
Eliminating the quadratic and cubic terms from these two equations
gives the displayed expression, proving necessity; substitution proves
sufficiency with exact coefficients. If $\Delta=0$, the sole possible
point is $-b/(3a)$ and it belongs exactly when $27a^2d=b^3$.
For degree two, $br^2+cr+d$ has the sole possible point $-c/(2b)$,
which belongs exactly when $c^2=4bd$. Degree zero or one gives none.
Finally omit $P=0$, which is outside ECC44. These finite algebraic
tests evaluate the full exceptional set without declaring that a
particular original period satisfies them.

Even when the quotient is bounded, the complete original state still
escapes. From FS13 and EQR5, $wq\to e_1/\lambda$, so
\[
 |t|^{1/2}\|\Phi_*q\|_U
   \longrightarrow\|\Phi_*e_1\|_U/|\lambda|>0,
 \qquad
 |t|^{1/2}\|\Psi_*q\|_W
   \longrightarrow\|\Psi_*e_1\|_W/|\lambda|>0.
 \tag{EQR12}
\]
If $U_1=0$, EQR3 returns its nonzero leading complement
$L_0M_*e_1/\lambda$. It cannot be zero, since EQR3 is invertible.
Equation EQR4 supplies its entire original norm and minimum
decomposition. Thus the quotient's exceptional cancellation has an
explicit nonzero location in the full original space.

\subsection{The cubic rational-sign escape cannot disappear in this quotient}
Now hold $t$, $P$ and the original receiver fixed, and use the original
affine variables $(a,y,z,w_0)$ with $y,z,w_0$ fixed as $a\to0$.
Here $a$ is the original ES coordinate, not $a_{\rm amp}$ and not
the collision constant $3P/22$. The full rational involution from
FS31 is
$(-a,y+2i/a,6i/a^2-z,w_0-14iy^2/a+28y/a^2+40i/a^3)$.
Applying the unchanged quotient gives exactly
\[
 \begin{aligned}
 E_tM_*\Sigma(a,y,z,w_0)
 ={}&40i a^{-3}u_4
 +a^{-2}(6iu_3+28y u_4)
 +a^{-1}(2iu_2-14iy^2u_4)\\
 &-a u_1+y u_2-z u_3+w_0u_4.
 \end{aligned}\tag{EQR13}
\]
Consequently its exact pole order is three if $u_4\ne0$, two if
$u_4=0,u_3\ne0$, and one if $u_4=u_3=0$.
In the last case $u_2\ne0$: otherwise the three independent columns
$M_*e_2,M_*e_3,M_*e_4$ would lie in the two-dimensional kernel of
$E_t$. The respective leading vectors are $40iu_4$, $6iu_3$ and
$2iu_2$, and their exact norm constants are their norms in
$G_{U,t}$ and $G_{W,t}$. This proves that no fixed two-root quotient
of the specified form hides the entire rational-sign escape.

These escaping vectors and the quotient operator have now both been
evaluated. OCQ11 still bounds every inverse exterior norm of the
induced original conductor; its scalar version uses each of the
two covariance Grams once. Divergence of an input vector in
EQR6 or EQR13 is compatible with this finite operator bound.
Equations EQR3--4 and EQR12 identify every lost direction and its
original norm, completing the correspondence without asserting an
original leading-moment zero or an RH endpoint.
